\documentclass[11pt,a4paper]{article}
\usepackage[utf8]{inputenc}
\usepackage[T1]{fontenc}
\usepackage{amsmath,amssymb,amsthm}
\usepackage{lmodern}
\usepackage[german]{babel}
\usepackage{geometry}
\geometry{margin=2.5cm}

\newtheorem{theorem}{Satz}[section]
\newtheorem{lemma}[theorem]{Lemma}
\newtheorem{definition}[theorem]{Definition}
\newtheorem{corollary}[theorem]{Korollar}

\title{Strukturelle Analyse des Beatty-Wandcodes und das Problem des permanenten Null-Lifts}
\author{Thomas Hoffbauer}
\date{\today}

\begin{document}

\maketitle

\begin{abstract}
Wir untersuchen die 2-adische und sturmianische Struktur des Beatty-Wandcodes, der durch die irrationale Steigung $\alpha = \log_2(9/4)$ gesteuert wird. Ziel ist der Ausschluss eines permanenten Null-Lift-Tails, woraus $x_* \notin \mathbb{N}_{>0}$ folgt. Wir präsentieren eine vollständige logische Architektur (\emph{Proof-Ladder}), die das Problem in drei bewiesene beziehungsweise direkt ableitbare Module (L1--L3) und einen einzigen verbleibenden globalen Kernsatz (L4: \emph{Past--Future Separation}) unterteilt.
\end{abstract}

\section{Einleitung und Problemstellung}

Der Gegenstand dieser Untersuchung ist die diophantische Analyse des Beatty-Wandcodes, welcher durch die Sprünge $s_K = \lfloor\alpha(K+1)\rfloor - \lfloor\alpha K\rfloor$ mit der Steigung $\alpha = \log_2(9/4)$ definiert ist. Das zentrale analytische Ziel besteht darin, rigoros zu zeigen, dass die Bahn des zugehörigen sturmianischen Pfades niemals in einen permanenten Null-Lift-Tail ($\ell_K = 0$ für alle $K \ge K_0$) „einfriert“. Ein solches Einfrieren würde äquivalent die Existenz eines ganzzahligen Grenzwerts $x_* \in \mathbb{N}$ bedeuten.

Mit Hilfe einer exakten 2-adischen Zustandsformulierung lässt sich das Problem auf die Frage reduzieren, ob der aus der diskreten Vergangenheit stammende rationale Präfixzustand $P_K$ jemals exakt den Wert der durch die irrationale Rotation $\theta_K = \{\alpha K\}$ gesteuerten Zukunftssumme $\Psi(\theta_K)$ annehmen kann.

\section{Die 2-adische Dynamik und das Reduktionsziel}

Aus der exakten Vorwärts-Rekurrenz für die rationalen Präfixe $P_K = 3^{2K} c_K$ und dem zugehörigen Null-Lift-Zukunftszustand $\Psi(\theta_K)$ definieren wir die Differenz:
$$D_K = P_K - \Psi(\theta_K).$$

Unter einem Null-Lift ($\ell_K = 0$) lautet die Rekurrenz:
$$D_{K+1} = \frac{9}{2^{2+s_K}} D_K.$$

In der 2-adischen Metrik besitzt der Faktor $\frac{9}{2^{2+s_K}}$ einen Betrag von $2^{2+s_K} > 1$. Die Null-Lift-Dynamik ist mithin \emph{expandierend}, nicht kontrahierend. Für die 2-adische Bewertung $\nu_K = v_2(D_K)$ gilt streng:
$$\nu_{K+1} = \nu_K - (2 + s_K).$$

Nach $N$ aufeinanderfolgenden Null-Lifts mit der kumulierten Match-Tiefe $L_{K,N}$ ergibt sich:
$$v_2(D_{K+N}) = v_2(D_K) - L_{K,N}.$$

Da $L_{K,N} \to \infty$ für $N \to \infty$, besitzt jedes von Null verschiedene $D_K$ nur eine endliche „Null-Lift-Reserve“. Ein permanenter Null-Tail ist folglich nur im exakten singulären Grenzfall möglich:
$$\boxed{D_K = 0 \iff \ell_j = 0 \quad \forall j \ge K \iff x_* = \rho_K \in \mathbb{N}.}$$

Damit ist das Problem auf den Nachweis reduziert, dass $D_K \neq 0$ für jedes $K \ge 0$ gilt.

\section{Die Beweisleiter: Etablierte Bausteine (L1--L3)}

Die Untersuchung gliedert sich in eine modulare Beweisleiter, deren erste drei Stufen vollständig abgesichert oder direkt herleitbar sind:

\subsection*{L1 --- Prefix Arithmetic}
Die exakte Kopplung der Präfixe und Liftziffern ist durch rationale Update-Gleichungen gegeben. Insbesondere lassen sich die lokalen Kongruenzbedingungen für die Liftziffer $\ell_K$ direkt auswerten:
$$s_K = 1 \implies \ell_K \equiv P_K - 3 \pmod 8,$$
$$s_K = 2 \implies \ell_K \equiv (9P_K + 1)3^{-2K-2} \pmod{16}.$$
Dies bildet das Fundament für die computergestützte Verifikation auf der Ebene kleiner Bit-Tiefen.

\subsection*{L2 --- Null-Lift Inverse Limit}
Für jedes $K$ und $N \ge 1$ existiert eine eindeutige Restklasse $r_{K,N} \pmod{2^{L_{K,N}}}$, welche die Bedingung für einen Null-Lift-Block der Länge $N$ codiert. Die geschachtelten Restklassen erfüllen:
$$r_{K,N+1} \equiv r_{K,N} \pmod{2^{L_{K,N}}}.$$
Der inverse Limes $\varprojlim_N r_{K,N} \in \mathbb{Z}_2$ existiert eindeutig und identifiziert sich exakt mit dem Zukunftspotential:
$$\boxed{r_{K,\infty} = \Psi(\theta_K).}$$

\subsection*{L3 --- First-Defect Valuation}
Wählen wir zwei Zukunftsphasen $\theta$ und $\theta'$ (beispielsweise entlang von Kettenbruch-Konvergenten $q_n$), deren sturmianische Wörter erstmals an einer Defektstelle $m$ abweichen, so lassen sich die Koeffizienten an dieser Stelle exakt bewerten. Die Differenz der ersten abweichend wirkenden Terme $\varepsilon - \varepsilon' = -\frac{1}{3} - \frac{1}{9} = -\frac{4}{9}$ liefert wegen $v_2(-4/9) = 2$ die exakte 2-adische Bewertung der Differenz:
$$\boxed{v_2\left(\Psi(\theta') - \Psi(\theta)\right) = L_m + 2.}$$
Dies ist keine asymptotische Abschätzung, sondern eine exakte diophantische Fixierung.

\section{Der offene Kern: Die Vergangenheit-Zukunft-Trennung (L4)}

Das gesamte verbleibende mathematische Problem ist in der vierten Stufe konzentriert:

\subsection*{L4 --- Past--Future Separation}
Zu zeigen ist der Hauptsatz:
$$\boxed{P_K \neq r_{K,\infty} \quad \forall K \ge 0.}$$
Äquivalent dazu in endlicher Form für Lean 4:
$$\boxed{\forall K \ge 0 \quad \exists N \ge 1: \quad P_K \not\equiv r_{K,N} \pmod{2^{L_{K,N}}}.}$$

Numerische Census-Daten für $K \le 5000$ zeigen, dass bereits eine Zertifikatstiefe von $M = 13$ ($P_K \bmod 8192$) in Kombination mit den drei elementaren Shifts $q \in \{1, 5, 6\}$ ausreicht, um über Bewertungsinkompatibilitäten ($v_2(Z_{K,q}) \neq T_{K,q}$) den Fall $D_K = 0$ für alle getesteten Präfixe auszuschließen. Der formale Nachweis, dass diese trennende Kongruenz universell für jedes beliebige $K$ existiert, ist der finale offene Schritt.

\section{Konditionale Folgerungen und Fazit}

Unter der Voraussetzung, dass L4 bewiesen ist, läuft die Widerlegung des permanenten Null-Tails mechanisch durch:
1. Aus L4 folgt $D_K = P_K - \Psi(\theta_K) \neq 0$.
2. Angenommen, es gäbe einen Null-Lift-Tail der Länge $N$ ab $K$, so liefert die dynamische Rekurrenz $v_2(D_{K+N}) = v_2(D_K) - L_{K,N}$.
3. Wegen $L_{K,N} \to \infty$ führt dies bei $D_K \neq 0$ zwingend zu einem Widerspruch zur 2-adischen Ganzzahligkeit, sofern der Tail unendlich fortbesteht.
4. Folglich müssen positive Liftziffern $\ell_j > 0$ unendlich oft auftreten, die Reste $\rho_K$ sind nicht schließlich konstant, und es gilt $x_* \notin \mathbb{N}_{>0}$.

Die Forschung ist damit auf einen einzigen, klar umrissenen Engpass reduziert: den Beweis der endlichen trennenden Kongruenz in L4.

\end{document}